
This work combines three research streams --- alignment bibliography, cross-community citation analysis, and S2AG bibliometric infrastructure --- none of which individually addresses directed citation asymmetry within an interdisciplinary alignment corpus.

\subsubsection{2.1 The huashen218 Alignment Corpus}

Shen et al. [2024] argue that AI alignment research is bidirectional: ML/NLP alignment (aligning AI behavior with human values) and HCI alignment (adapting humans to AI systems) are distinct but mutually dependent. Their systematic review assembled the huashen218 corpus --- over 400 papers from NeurIPS, ICML, ICLR, ACL, EMNLP, CHI, CSCW, IUI, and adjacent venues --- as the operational definition of the bidirectional alignment community. The ICLR 2025 Workshop on Bidirectional Human-AI Alignment further validated this framing, identifying ML/NLP and HCI as the primary disciplinary clusters.

Shen et al. [2024] characterize the community through paper authorship and venue affiliation, not through citation structure. No prior work has constructed a directed citation graph within the huashen218 corpus or empirically tested whether the claimed bidirectionality manifests in how papers from the two communities cite each other.

\subsubsection{2.2 Cross-Community Citation Asymmetry}

The closest methodological precedent is Wahle et al. [2023], who analyzed cross-field citation influence in NLP using S2AG bulk data. Their analysis of 77,000 NLP papers found strong within-field citation preference and measurable cross-field asymmetry, validating the use of S2AG directed citation edges and $2\times 2$ contingency tables for cross-community citation analysis. The same statistical framework is adopted here --- chi-squared test on directed citation proportions --- applied to an interdisciplinary alignment-specific corpus rather than a single-field corpus.

Chen [2023] analyzed HCI community self-citation rates using the X-index metric across CHI, UIST, and CSCW papers from 2010--2020, finding HCI self-citation rates increasing over time. The present finding that HCI alignment papers in the huashen218 proxy show zero $HCI\rightarrow HCI$ within-corpus citations is complementary but distinct: alignment-specific HCI papers appear to cite ML/NLP alignment work exclusively as within-corpus outgoing references, which is consistent with cross-field dependency rather than self-referential community growth.

\subsubsection{2.3 Bibliometric Classification Methods}

Standard bibliometric classification relies on venue name matching. This approach is adequate for single-community corpora where venue names appear in abbreviated form. However, S2AG returns full proceedings names in the \texttt{venue} field, making substring matching against abbreviated venue strings unreliable for multi-community corpora. The S2AG API \citep{Wade2022S2AG} exposes the \texttt{fieldsOfStudy} field as a tag-based disciplinary classifier derived from S2AG's knowledge graph, populated independently of venue name formatting. The present work formalizes FoS-primary classification as a distinct scheme, reports the 100\% versus 12\% coverage advantage, and documents the root cause to support replication.

\subsubsection{2.4 Positioning}

This work takes Shen et al. [2024]'s corpus as input, applies Wahle et al. [2023]'s statistical framework to it, and identifies that FoS-primary classification is required to make this combination function on an interdisciplinary corpus. The methodological gap --- robust classification for interdisciplinary alignment corpora using S2AG data --- is the primary methodological contribution, alongside the first empirical citation asymmetry measurement in the huashen218 corpus.

